% =============================================================================
\chapter{O hardware}
\label{cap:hardware}
% =============================================================================

\section{A placa DE1-SoC}

A DE1-SoC, da Terasic, é construída em torno de um Cyclone~V SoC da Intel
(\texttt{5CSEMA5F31C6}). O ``SoC'' do nome quer dizer que o chip reúne duas
coisas bem diferentes:

\begin{description}
  \item[HPS] (\emph{Hard Processor System}): um processador ARM Cortex-A9 de
    dois núcleos, com memória e periféricos próprios, que roda um Linux gravado
    no cartão SD da placa. No Morphe ele cuida da rede e roda o servidor.
  \item[FPGA] lógica programável, configurada pelo bitstream que o Quartus gera.
    No Morphe ela contém os aceleradores e as memórias em que eles trabalham.
\end{description}

Os dois lados estão no mesmo chip e se comunicam por pontes internas, sem
passar por nenhum fio da placa. É essa proximidade que permite ao servidor
escrever mil amostras numa memória da FPGA como se escrevesse num vetor em C.

A lógica da FPGA roda com o relógio de 50\,MHz da placa (\texttt{CLOCK\_50}).
A análise de tempo da última compilação (15/09/2026) mede uma frequência máxima
de 60,99\,MHz para esse domínio: há folga de cerca de 20\,\%. A
Tabela~\ref{tab:recursos} mostra quanto do chip o projeto ocupa.

\begin{table}[htbp]
  \centering
  \small
  \caption{Recursos da FPGA ocupados pelo projeto (compilação de 15/09/2026).}
  \label{tab:recursos}
  \begin{tabular}{@{}lrrr@{}}
    \toprule
    \textbf{Recurso} & \textbf{Usado} & \textbf{Disponível} & \textbf{Ocupação} \\
    \midrule
    Elementos lógicos (ALMs)      & 11\,140 & 32\,070 & 35\,\% \\
    Blocos DSP (multiplicadores)  & 30      & 87      & 34\,\% \\
    Blocos de memória M10K        & 165     & 397     & 42\,\% \\
    \bottomrule
  \end{tabular}
\end{table}

O que sobra é bastante: cerca de 230 blocos M10K livres, pouco menos de
300\,KB, o suficiente para uma memória de captura de dezenas de milhares de
amostras quando a aquisição pelo ADC for implementada (seção~\ref{sec:adc}).

\subsection*{O bitstream e a licença da FFT}

O bitstream versionado é \arq{Quartus/output_files/soc_system_time_limited.sof},
compilado com o Quartus Prime Lite 20.1. O núcleo de FFT é um IP da Intel que a
edição Lite só permite usar em modo de avaliação: o bitstream funciona
\textbf{enquanto o programador JTAG continuar conectado} à placa --- o
\emph{tether} --- e para depois de cerca de uma hora sem ele. Isso vale para o
bitstream inteiro, não só para a FFT: sem o \emph{tether}, a convolução também
passa a esgotar o tempo. O \arq{morphe-up.sh} mantém o \emph{tether} vivo
(capítulo~\ref{cap:laboratorio}).

\section{Pontes e mapa de memória}
\label{sec:pontes}

O HPS enxerga a FPGA como faixas de endereços. O Morphe usa duas pontes
(Figura~\ref{fig:mapa-memoria}):

\begin{description}
  \item[Ponte HPS\(\to\)FPGA] (\texttt{h2f\_axi\_master}), a partir de
    \texttt{0xC000\,0000}. Larga e rápida, é por onde passam os dados: as
    memórias de entrada e de saída de cada acelerador.
  \item[Ponte leve] (\emph{lightweight}), a partir de \texttt{0xFF20\,0000}.
    Estreita, feita para registradores de controle: os PIOs de
    \texttt{start}, \texttt{done} e de parâmetros.
\end{description}

\begin{figure}[htbp]
  \centering
  % Mapa de memoria. Enderecos de C/hps_0.h e C/address_map_arm.h.
\begin{tikzpicture}[
    reg/.style={draw, minimum width=34mm, minimum height=5.2mm, inner sep=1pt,
                font=\scriptsize\ttfamily, anchor=north},
    end/.style={font=\scriptsize\ttfamily, anchor=east, text=black!70},
  ]
  % ---- Ponte HPS->FPGA ---------------------------------------------------
  \node[font=\small\bfseries, align=center] at (0,8mm) {Ponte HPS\(\to\)FPGA\\\normalfont\scriptsize base \texttt{0xC000\,0000}, 32 bits};
  \foreach \i/\nome/\off/\cor in {
      0/fir\_yn/0x10000/opConv,
      1/conv1d\_yn/0x12000/opConv,
      2/fir\_hn/0x14000/opConv,
      3/fir\_xn/0x15000/opConv,
      4/conv1d\_hn/0x16000/opConv,
      5/conv1d\_xn/0x17000/opConv,
      6/fft\_yn\_re/0x18000/opFft,
      7/fft\_yn\_imag/0x19000/opFft,
      8/fft\_xn\_re/0x1A000/opFft,
      9/fft\_xn\_imag/0x1B000/opFft,
      10/iir\_xn/0x1C000/opIir,
      11/iir\_yn/0x1D000/opIir,
      12/iir\_coef/0x1E000/opIir}
  {
    \node[reg, fill=\cor!12] (h\i) at (0,-\i*5.2mm) {\nome};
    \node[end] at (h\i.west) {\off\ };
  }

  % ---- Ponte leve ---------------------------------------------------------
  \begin{scope}[xshift=72mm]
  \node[font=\small\bfseries, align=center] at (0,8mm) {Ponte leve (LW)\\\normalfont\scriptsize base \texttt{0xFF20\,0000}, PIOs};
  \foreach \i/\nome/\off/\cor in {
      0/fir\_error/0x00/opConv,
      1/fir\_done/0x10/opConv,
      2/fir\_start/0x20/opConv,
      3/conv1d\_done/0x30/opConv,
      4/conv1d\_start/0x40/opConv,
      5/fft\_bfp\_exponent/0x50/opFft,
      6/fft\_wrapper\_done/0x60/opFft,
      7/fft\_inverse/0x70/opFft,
      8/fft\_wrapper\_start/0x80/opFft,
      9/iir\_start/0x90/opIir,
      10/iir\_done/0xA0/opIir,
      11/iir\_error/0xB0/opIir,
      12/iir\_nsecoes/0xC0/opIir}
  {
    \node[reg, fill=\cor!12] (l\i) at (0,-\i*5.2mm) {\nome};
    \node[end] at (l\i.west) {\off\ };
  }
  \end{scope}

  % Legenda
  \begin{scope}[yshift=-72mm, xshift=-10mm]
    \foreach \x/\c/\t in {0/opConv/convolução e FIR, 38/opFft/FFT e IFFT, 72/opIir/IIR}
      { \fill[\c!12, draw=black] (\x mm,0) rectangle ++(4mm,3mm);
        \node[anchor=west, font=\scriptsize] at (\x mm+5mm,1.5mm) {\t}; }
  \end{scope}
\end{tikzpicture}

  \caption{Mapa de memória visto pelo servidor: memórias de dados atrás da
    ponte HPS\(\to\)FPGA e registradores de controle atrás da ponte leve.
    Os endereços são deslocamentos a partir da base de cada ponte.}
  \label{fig:mapa-memoria}
\end{figure}

\paragraph{As memórias têm duas portas.} Cada uma das treze memórias de dados
é uma RAM interna da FPGA de porta dupla: uma porta está ligada à ponte, para o
HPS, e a outra ao acelerador. O servidor e o acelerador nunca disputam a mesma
porta; o que os separa no tempo é o protocolo de \texttt{start} e \texttt{done}
(seção~\ref{sec:aceleradores}). As memórias de entrada guardam 1024 palavras de
32 bits (4\,KB); as de saída da convolução e do FIR, 2047 (8\,KB); a de
coeficientes do IIR, 5 palavras por seção para até 16 seções.

\paragraph{Como o servidor chega lá.} O servidor abre \arq{/dev/mem} e mapeia
as duas faixas no próprio espaço de endereços com \texttt{mmap}; a partir daí,
escrever \texttt{g\_conv\_xn[i]} é escrever na memória da FPGA. Por isso ele
precisa rodar como \texttt{root}. O tamanho da faixa mapeada é calculado em
tempo de compilação a partir do \arq{hps_0.h}, e verificações estáticas no
código recusam compilar se alguma memória for menor do que os limites de
\arq{morphe_config.h} exigem.

\paragraph{O \arq{hps_0.h} é gerado, nunca editado.} Os endereços da
Figura~\ref{fig:mapa-memoria} são decididos no Platform Designer, a ferramenta
do Quartus que monta o sistema, e ficam registrados no arquivo
\arq{soc_system.sopcinfo}. O \arq{Quartus/gen_hps_header.py} lê esse arquivo e
escreve o \arq{C/hps_0.h}; com \opt{check}, apenas confere se o cabeçalho
versionado ainda corresponde ao projeto. Um servidor compilado com endereços
velhos escreve nos lugares errados sem dar erro nenhum; a conferência existe
para que isso não aconteça.

\medskip
\noindent\fbox{\parbox{\dimexpr\textwidth-2\fboxsep-2\fboxrule}{\small
\textbf{Cuidado com o mapa do ``DE1-SoC Computer''.} Os exemplos da Intel
para a placa usam \texttt{0xC800\,0000} como base da memória da FPGA. No
Morphe a base é \texttt{0xC000\,0000}. Código copiado desses exemplos lê e
escreve fora das memórias do projeto (\arq{docs/RESSALVAS.md}, item 11).}}

\section{Os aceleradores}
\label{sec:aceleradores}

A Figura~\ref{fig:fpga-topo} mostra o que há dentro da FPGA: quatro
aceleradores, cada um com as suas memórias e os seus registradores de
controle, todos ligados ao HPS pela interconexão que o Platform Designer gera.

\begin{figure}[htbp]
  \centering
  % Topo da FPGA: cada acelerador com as suas memorias e PIOs.
% Linha por acelerador: [SRAMs de entrada] -> [nucleo] -> [SRAM de saida].
\begin{tikzpicture}[
    ent/.style={mem, minimum width=17mm},
    nuc/.style={caixa, fill=white, minimum width=34mm, minimum height=13mm, font=\scriptsize},
    pio/.style={font=\tiny\ttfamily, text=black!70},
  ]
  \def\xe{0}\def\xn{36mm}\def\xs{72mm}

  % --- convolucao -------------------------------------------------------
  \node[ent, fill=opConv!10] (cx) at (\xe,  3mm) {conv1d\_xn};
  \node[ent, fill=opConv!10] (ch) at (\xe, -3mm) {conv1d\_hn};
  \node[nuc, draw=opConv]    (cc) at (\xn, 0)    {\texttt{conv1d}\\MAC Q15.16\\\(x, h \le 1024 \to y \le 2047\)};
  \node[ent, fill=opConv!10] (cy) at (\xs, 0)    {conv1d\_yn};

  % --- FIR ----------------------------------------------------------------
  \node[ent, fill=opConv!10] (fx) at (\xe, -19mm) {fir\_xn};
  \node[ent, fill=opConv!10] (fh) at (\xe, -25mm) {fir\_hn};
  \node[nuc, draw=opConv]    (fc) at (\xn, -22mm) {\texttt{conv1d} (\texttt{fir\_inst})\\o mesmo RTL da convolução};
  \node[ent, fill=opConv!10] (fy) at (\xs, -22mm) {fir\_yn};

  % --- FFT ----------------------------------------------------------------
  \node[ent, fill=opFft!10] (tx) at (\xe, -41mm) {fft\_xn\_re};
  \node[ent, fill=opFft!10] (ti) at (\xe, -47mm) {fft\_xn\_imag};
  \node[nuc, draw=opFft]    (tc) at (\xn, -44mm) {\texttt{fft\_wrapper} + IP FFT II\\Buffered Burst, N = 1024\\Q15.8, expoente de bloco};
  \node[ent, fill=opFft!10] (ty) at (\xs, -41mm) {fft\_yn\_re};
  \node[ent, fill=opFft!10] (tj) at (\xs, -47mm) {fft\_yn\_imag};

  % --- IIR ----------------------------------------------------------------
  \node[ent, fill=opIir!10] (ix) at (\xe, -63mm) {iir\_xn};
  \node[ent, fill=opIir!10] (ik) at (\xe, -69mm) {iir\_coef};
  \node[nuc, draw=opIir]    (ic) at (\xn, -66mm) {\texttt{iir\_cascade}\\até 16 seções, Q15.16\\1 MAC (15 blocos DSP)};
  \node[ent, fill=opIir!10] (iy) at (\xs, -66mm) {iir\_yn};

  % Setas de dados
  \foreach \a/\b in {cx/cc,ch/cc,fx/fc,fh/fc,tx/tc,ti/tc,ix/ic,ik/ic}
    \draw[seta] (\a.east) -- (\a.east -| \b.west);
  \foreach \a/\b in {cc/cy,fc/fy,ic/iy}
    \draw[seta] (\a.east) -- (\b.west);
  \draw[seta] (tc.east |- ty) -- (ty.west);
  \draw[seta] (tc.east |- tj) -- (tj.west);

  % PIOs de controle (acima de cada nucleo)
  \node[pio, above=0.5mm of cc] {PIO: start, done};
  \node[pio, above=0.5mm of fc] {PIO: start, done, error};
  \node[pio, above=0.5mm of tc] {PIO: start, done, inverse, bfp\_exponent};
  \node[pio, above=0.5mm of ic] {PIO: start, done, error, nsecoes};

  % Interconexao: desce pela esquerda (entradas) e pela direita (saidas)
  \coordinate (bE) at (-12mm, 0);        % barramento da esquerda
  \coordinate (bD) at (\xs+13mm, 0);     % barramento da direita
  \def\ytopo{14mm}
  \draw[thick, cHps] (bE |- 0,-72mm) -- (bE |- 0,\ytopo) -- node[rotulo, above, text=cHps]{interconexão do Platform Designer} (bD |- 0,\ytopo) -- (bD |- 0,-66mm);
  \foreach \y in {3,-3,-19,-25,-41,-47,-63,-69}
    \draw[cHps] (bE |- 0,\y mm) -- (\xe-8.5mm,\y mm);
  \foreach \y in {0,-22,-41,-47,-66}
    \draw[cHps] (bD |- 0,\y mm) -- (\xs+8.5mm,\y mm);

  % HPS, a esquerda, fora da FPGA
  \node[hps, minimum height=40mm, text width=21mm, anchor=east] (hps) at (-24mm,-29mm)
    {HPS\\[2mm]\scriptsize ponte HPS\(\to\)FPGA: memórias\\[2mm] ponte leve: PIOs};
  \draw[<->, very thick, cHps] (hps.east) -- (bE |- hps.east);

  % cantos da FPGA (virgula dentro de fit=... quebraria a lista de opcoes)
  \coordinate (cantoA) at (bE |- 0,\ytopo);
  \coordinate (cantoB) at (bD |- 0,-72mm);
  \begin{scope}[on background layer]
    \node[grupo, draw=cFpga, fill=cFpga!4, fit=(cx)(iy)(ik)(tj)(cantoA)(cantoB),
          inner ysep=5mm,
          label={[font=\scriptsize\bfseries, text=cFpga]above:{FPGA (Cyclone V), relógio de 50 MHz}}] {};
  \end{scope}
\end{tikzpicture}

  \caption{O que há dentro da FPGA: quatro aceleradores, cada um com as suas
    memórias e os seus registradores de controle.}
  \label{fig:fpga-topo}
\end{figure}

\subsection{O protocolo comum}

Todos os aceleradores obedecem ao mesmo combinado com o servidor:

\begin{enumerate}
  \item com \texttt{start} em 0, o servidor escreve as entradas nas memórias;
  \item o servidor leva \texttt{start} a 1; o acelerador detecta a
    \emph{subida} do sinal, zera o \texttt{done} e começa;
  \item o acelerador lê as entradas, calcula, escreve as saídas e leva
    \texttt{done} a 1;
  \item o servidor, que lê \texttt{done} a cada 50\,\textmu s, lê as saídas e devolve
    \texttt{start} a 0.
\end{enumerate}

Como o acelerador reage à subida e não ao nível, um \texttt{start} esquecido em
1 não dispara uma segunda execução. E como o servidor zera todos os
\texttt{start} na primeira requisição depois de a FPGA ser preparada, um valor
que tenha sobrado de antes do bitstream do Morphe também não dispara nada.

Os aceleradores acessam as memórias por pequenos controladores que gastam
alguns ciclos por acesso. É um custo deliberado de simplicidade: numa
convolução, o tempo é dominado por esses acessos, não pela multiplicação. Os
controladores estão em \arq{memory_read_controller.v} e em
\arq{memory_write_controller.v}.

\paragraph{O tempo como sinal de saúde.} Uma convolução de 1024 por 1024
amostras executa cerca de 8,4 milhões de ciclos, uns 168\,ms a 50\,MHz; de ponta
a ponta, com a rede, o cliente mede 214\,ms. Uma resposta muito mais rápida que
isso é sinal de defeito, não de sorte: um bitstream mal compilado já foi visto
levantar o \texttt{done} em 3,3\,ms sem ter calculado nada
(\arq{docs/RESSALVAS.md}, item 13).

\subsection{Convolução e FIR: o \texttt{conv1d}}
\label{sec:acel-conv}

O \arq{conv1d.v} calcula a convolução linear
\[
  y[n] = \sum_{k} x[k]\,h[n-k], \qquad n = 0, 1, \dots, 2046,
\]
com os limites da soma ajustados a cada \(n\) para não ler fora dos vetores. É
uma máquina de estados com um único multiplicador-acumulador: para cada \(n\),
lê \(x[k]\) e \(h[n-k]\), multiplica, acumula e, ao fim da soma, escreve
\(y[n]\).

A aritmética é em Q15.16. O produto de dois números de 32 bits tem 64 bits;
o \texttt{conv1d} guarda os bits 47 a 16, o que realinha o ponto fracionário
\emph{truncando} os 16 bits de baixo. O acumulador tem 32 bits e \emph{não
satura}: quando a soma passa de \(\pm 32768\) ela dá a volta e aparece com o
sinal trocado. Duas consequências práticas:

\begin{itemize}
  \item o erro de truncamento se acumula com o número de termos: numa
    convolução de 1024 por 1024 ele chega a centenas de LSB. É esperado;
  \item com sinais de amplitude \(A\) e \(N\) termos, a saída estoura quando
    \(N A^2 \ge 32768\) --- com \(N = 1024\), a partir de \(A \approx 5{,}7\).
    Uma saída negativa para sinais positivos é esse estouro, não um defeito.
\end{itemize}

Trocar o truncamento por arredondamento reduziria o erro máximo medido de 521
para 25 LSB. A troca foi avaliada e \textbf{deliberadamente não feita}
(14/09/2026): o comportamento numérico do \texttt{conv1d} é o que o TCC
original descreve, e ele fica registrado em \arq{docs/RESSALVAS.md} como
escolha, não como defeito esquecido.

\paragraph{O FIR é o mesmo circuito.} O filtro FIR é uma segunda instância do
\texttt{conv1d} (\texttt{fir\_inst}), com memórias próprias. Filtrar um sinal
por um FIR é convoluí-lo com a resposta ao impulso do filtro, e é exatamente
isso que acontece; a diferença entre as duas operações está no aplicativo, que
oferece para o FIR um projetista de filtros. O projeto ainda contém um IP de FIR
da Intel (\texttt{fir\_ii}) e o seu adaptador, \arq{fir_wrapper.v}, mas a
instância está comentada no \arq{ghrd_top.v}. O registrador
\texttt{fir\_error} sobrou desse adaptador e o \texttt{conv1d} não o aciona.

\subsection{FFT e IFFT: o núcleo da Intel}
\label{sec:acel-fft}

A transformada usa o IP \emph{FFT~II} da Intel, configurado para 1024 pontos
em modo \emph{Buffered Burst}, com entrada de 24 bits e fatores de giro de
18 bits. O \arq{fft_wrapper.v} faz a ponte entre as memórias e o IP: lê as
partes real e imaginária da entrada, entrega-as ao IP amostra por amostra
respeitando o \emph{handshake} do barramento de fluxo, recebe o resultado e o
escreve nas memórias de saída. O registrador \texttt{fft\_inverse} escolhe entre
a transformada direta e a inversa; o resto do caminho é o mesmo.

Duas particularidades desse núcleo aparecem no resultado:

\begin{description}
  \item[Expoente de bloco.] A cada estágio, o IP detecta se os valores
    cresceram e, se preciso, desloca as 1024 amostras \emph{juntas},
    acumulando um expoente único para o bloco (\emph{block floating point}). O
    adaptador captura esse expoente no registrador \texttt{fft\_bfp\_exponent},
    e o servidor o aplica, junto com o fator \(2^{-8}\) do formato Q15.8, antes
    de devolver o resultado em ponto flutuante. É a única conversão numérica
    feita no servidor.
  \item[IFFT sem o fator \(1/N\).] O núcleo entrega a inversa não normalizada,
    \(N\) vezes a IDFT matemática. O cliente aplica o \(1/1024\).
\end{description}

O capítulo~\ref{cap:numerica} explica o que o expoente de bloco e os 24 bits
de entrada significam para a precisão, e por que ela para em torno de 92\,dB.

\subsection{IIR: a cascata de seções de segunda ordem}
\label{sec:acel-iir}

O filtro IIR é a única operação em que a saída depende das saídas anteriores,
e isso muda o que importa no projeto do circuito. Ele é implementado como uma
cascata de até 16 seções de segunda ordem (\emph{biquads}), cada uma em forma
direta~I (Figura~\ref{fig:iir-cascata}):
\[
  y_s[n] = b_0\,x_s[n] + b_1\,x_s[n-1] + b_2\,x_s[n-2]
           - a_1\,y_s[n-1] - a_2\,y_s[n-2],
\]
com a saída de cada seção servindo de entrada para a seguinte. Os coeficientes
vêm da memória \texttt{iir\_coef}, cinco por seção, na ordem
\(b_0, b_1, b_2, a_1, a_2\) (\(a_0 = 1\) fica implícito); o número de seções,
do registrador \texttt{iir\_nsecoes}.

\begin{figure}[htbp]
  \centering
  % Cascata de secoes de 2a ordem e uma secao em forma direta I.
\begin{tikzpicture}[
    sos/.style={caixa, draw=opIir, fill=opIir!10, minimum width=16mm},
    som/.style={circle, draw, inner sep=0pt, minimum size=5mm},
    atr/.style={draw, minimum size=6mm, font=\scriptsize, inner sep=1pt},
    gan/.style={isosceles triangle, draw, inner sep=0.6pt, minimum height=5mm,
                font=\tiny, isosceles triangle apex angle=60},
  ]
  % ---- Cascata -----------------------------------------------------------
  \node (xin) at (0,0) {\(x[n]\)};
  \node[sos, right=8mm of xin] (s1) {seção 1};
  \node[sos, right=8mm of s1]  (s2) {seção 2};
  \node[right=6mm of s2] (dots) {\(\cdots\)};
  \node[sos, right=6mm of dots] (sS) {seção \(S\)};
  \node[right=8mm of sS] (yout) {\(y[n]\)};
  \draw[seta] (xin) -- (s1);
  \draw[seta] (s1) -- (s2);
  \draw[seta] (s2) -- (dots);
  \draw[seta] (dots) -- (sS);
  \draw[seta] (sS) -- (yout);
  \node[rotulo, below=1mm of s2] {cada seção: \(b_0, b_1, b_2, a_1, a_2\) em Q15.16, com \(a_0 = 1\)\\
     até 16 seções, ordem até 32};

  % ---- Uma secao, forma direta I -------------------------------------------
  \begin{scope}[yshift=-26mm, xshift=10mm]
    \node (x) at (0,0) {\(x_s[n]\)};
    \coordinate (xa) at (12mm,0);
    \node[atr] (z1) at (12mm,-11mm) {\(z^{-1}\)};
    \node[atr] (z2) at (12mm,-22mm) {\(z^{-1}\)};
    \node[gan] (b0) at (22mm,0)    {\(b_0\)};
    \node[gan] (b1) at (22mm,-11mm){\(b_1\)};
    \node[gan] (b2) at (22mm,-22mm){\(b_2\)};
    \node[som] (sm) at (38mm,0) {\(+\)};
    \node[caixa, fill=white, font=\scriptsize] (q) at (56mm,0) {arredonda\\satura};
    \coordinate (ya) at (72mm,0);
    \node (y) at (82mm,0) {\(y_s[n]\)};
    \node[atr] (w1) at (72mm,-11mm) {\(z^{-1}\)};
    \node[atr] (w2) at (72mm,-22mm) {\(z^{-1}\)};
    \node[gan, shape border rotate=180] (a1) at (54mm,-11mm) {\(-a_1\)};
    \node[gan, shape border rotate=180] (a2) at (54mm,-22mm) {\(-a_2\)};

    \draw[seta] (x) -- (b0);
    \draw[seta] (xa) -- (z1);  \draw[seta] (z1) -- (z2);
    \draw[seta] (z1) -- (b1);  \draw[seta] (z2) -- (b2);
    \draw[seta] (b0) -- (sm);
    \draw[seta] (b1) -| (sm);
    \draw[seta] (b2) -| ($(sm.south)+(-1.5mm,0)$);
    \draw[seta] (sm) -- (q);
    \draw[seta] (q) -- (y);
    \draw[seta] (ya) -- (w1);  \draw[seta] (w1) -- (w2);
    \draw[seta] (w1) -- (a1);  \draw[seta] (w2) -- (a2);
    \draw[seta] (a1) -| ($(sm.south)+(1.5mm,0)$);
    \draw[seta] (a2) -| ($(sm.south)+(1.5mm,0)$);
    \fill (xa) circle (0.6mm);  \fill (ya) circle (0.6mm);

    \node[rotulo, align=left, anchor=west] at (86mm,-14mm)
      {4 registradores\\de estado por seção;\\as 5 multiplicações\\num só MAC\\compartilhado};
  \end{scope}
\end{tikzpicture}

  \caption{O bloco IIR: cascata de seções de segunda ordem (acima) e a seção em
    forma direta~I (abaixo). Um só multiplicador-acumulador atende todas as
    seções, dois ciclos por seção.}
  \label{fig:iir-cascata}
\end{figure}

Cada escolha desse projeto foi medida antes de ser feita:

\begin{description}
  \item[Seções de segunda ordem, e não um polinômio de ordem \(N\).] A posição
    dos polos de um polinômio de ordem alta é extremamente sensível aos
    coeficientes. Um Butterworth de ordem 22, estável por construção, volta do
    próprio MATLAB com um polo de módulo 1,278 --- instável --- porque achar as
    raízes daquele polinômio já é mal condicionado em ponto flutuante de 64
    bits. A mesma especificação em seções de segunda ordem tem o maior polo em
    0,988.
  \item[Forma direta I.] Tem um acumulador único e largo e arredonda uma vez
    só, na saída da seção. A forma direta II transposta, melhor em ponto
    flutuante, guarda estados já arredondados, e numa realimentação esse erro
    circula.
  \item[Arredondar, e não truncar.] Truncar introduz um viés de meio LSB
    negativo por amostra. Num filtro realimentado esse viés não some: vira um
    deslocamento DC de \(-0{,}5/(1 + a_1 + a_2)\) LSB na saída. No FIR não
    importa, e por isso o \texttt{conv1d} pôde manter o truncamento.
  \item[Q15.16, o mesmo formato do \texttt{conv1d}.] É suficiente: o raio dos
    polos se desloca no máximo cerca de \(10^{-6}\) ao quantizar, mesmo com
    polos em \(|p| = 0{,}9984\).
  \item[Uma amostra atravessa todas as seções antes da próxima.] Assim bastam
    quatro registradores de estado por seção e uma única passagem pela
    memória. O resultado é idêntico, bit a bit, ao de processar seção por
    seção, porque cada seção só vê a saída já arredondada da anterior.
  \item[Um só multiplicador-acumulador, compartilhado.] As cinco multiplicações
    de 32 por 32 bits de uma seção são feitas juntas, e o mesmo circuito
    atende as 16 seções em sequência. O custo em blocos DSP (15, porque cada
    multiplicação de 32 bits ocupa três multiplicadores de 27 bits) não depende
    do número de seções. Cada seção leva dois ciclos: um para selecionar os
    operandos, outro para calcular. Com um ciclo só, o circuito não fechava a
    50\,MHz (49,11\,MHz medidos em 14/09/2026); com dois, fecha a 60,99\,MHz.
\end{description}

Se alguma amostra satura, o registrador \texttt{iir\_error} sobe; o resultado
é devolvido mesmo assim, saturado, e o cliente avisa o aluno. O estado interno é
zerado a cada execução, para que o mesmo sinal dê sempre a mesma saída.

\paragraph{Verificado bit a bit.} O RTL do IIR tem um modelo de referência em
Python que usa inteiros de precisão arbitrária (o acumulador de uma seção
chega a \(2^{65}\), além do que um \texttt{float64} representa exatamente). O
testbench em Icarus Verilog (\arq{Quartus/roda_tb_iir.sh}) e a placa
(\arq{testa_iir_hw.py}) são comparados com esse modelo e precisam coincidir
em todas as amostras: num filtro realimentado, divergir um LSB numa amostra é
divergir para sempre dali em diante, então ``parecido'' não serve de critério.

\section{Periféricos ainda não usados: ADC e codec}
\label{sec:adc}

A placa tem dois conversores que o Morphe ainda não usa. Esta seção descreve o
estado de hoje e o que falta; a Figura~\ref{fig:adc} resume as duas coisas.

\begin{figure}[htbp]
  \centering
  % ADC: o que existe hoje (linha cheia) e o que falta (tracejado).
\begin{tikzpicture}[node distance=6mm]
  % ---- Hoje ---------------------------------------------------------------
  \node[caixa, fill=white, text width=16mm] (con) {conector\\analógico\\\scriptsize CH0--CH7};
  \node[caixa, fill=black!6, right=of con, text width=23mm] (ltc) {LTC2308\\\scriptsize 12 bits, 8 canais\\0--4,095\,V};
  \node[fpga, right=18mm of ltc, text width=31mm] (up) {controlador do\\University Program\\\scriptsize\itshape instância comentada};

  \draw[seta] (con) -- (ltc);
  \draw[<->, thick] (ltc) -- node[rotulo, above]{SPI, 4 fios}
                              node[rotulo, below]{\texttt{CONVST}\\\texttt{SCLK DIN}\\\texttt{DOUT}} (up);

  % ---- O que falta ----------------------------------------------------------
  \node[futuro, right=8mm of up, text width=18mm] (pio) {PIO\\\scriptsize leitura avulsa};
  \node[futuro, below=14mm of ltc, text width=24mm] (ctl) {controlador próprio\\\scriptsize CONVST a \(1/f_s\)};
  \node[futuro, text width=18mm] (ram) at (ctl -| up) {RAM de\\captura};
  \node[futuro, text width=18mm] (op) at (ctl -| pio) {opcodes;\\janela\\``Aquisição''};

  \draw[seta, dashed, cFuturo] (up) -- (pio);
  \draw[seta, dashed, cFuturo] (pio) -- (op);
  \draw[seta, dashed, cFuturo] (ltc) -- (ctl);
  \draw[seta, dashed, cFuturo] (ctl) -- (ram);
  \draw[seta, dashed, cFuturo] (ram) -- (op);

  \node[rotulo, anchor=north, text=cFuturo!70!black, align=left] at ($(ctl.south)!0.5!(op.south)+(0,-2mm)$)
    {etapa 1: leitura avulsa, o \emph{voltímetro} (linha de cima)\\
     etapa 2: captura a \(f_s\) fixa em RAM (linha de baixo)};
\end{tikzpicture}

  \caption{O ADC da placa: o que existe (linha cheia) e o que falta para
    capturar sinais a uma taxa fixa (tracejado).}
  \label{fig:adc}
\end{figure}

\paragraph{O ADC.} O conversor analógico-digital da DE1-SoC é o LTC2308: 12
bits, 8 canais, por aproximações sucessivas, ligado à FPGA por uma interface
serial de quatro fios (\texttt{ADC\_CONVST}, \texttt{ADC\_SCLK},
\texttt{ADC\_DIN}, \texttt{ADC\_DOUT}). As entradas analógicas chegam por um
conector próprio da placa. O projeto contém o controlador do \emph{University
Program} da Intel para esse conversor, configurado para entrada simples e
unipolar: com a referência interna, a faixa é de 0 a 4,095\,V, 1\,mV por passo.
\textbf{Tensão negativa não é lida e pode danificar a entrada}; um sinal
alternado precisa de um deslocamento DC antes de chegar ao conector.

Hoje a instância desse controlador está comentada no \arq{ghrd_top.v}: o ADC
não faz parte do bitstream. E mesmo ativo, ele não serviria para amostrar um
sinal: converte sem parar, sozinho, e entrega apenas o último valor lido de
cada canal, sem controle do instante da conversão. Uma incerteza de alguns
microssegundos no instante da amostra limita a medida de um tom de 5\,kHz a
cerca de 25\,dB. Capturar a uma taxa de amostragem fixa exige um controlador
próprio, que dispare a conversão a cada \(1/f_s\) e guarde as amostras numa
memória da FPGA. O plano, em duas etapas --- uma leitura avulsa, como um
voltímetro, e depois a captura temporizada --- está em \arq{docs/ADC.md}.

\paragraph{O codec de áudio.} A placa tem também um codec de áudio (o WM8731,
segundo a documentação da DE1-SoC), com entrada e saída de linha, ligado à FPGA
por uma interface de áudio serial e configurado por I\textsuperscript{2}C. Os
pinos estão declarados no \arq{ghrd_top.v} e mais nada: nenhum sinal do codec é
usado pelo projeto. Seria o caminho natural para uma saída de áudio --- ouvir o
sinal filtrado --- e fica como trabalho futuro.
